% Compile from this directory:
%   pdflatex financial_sde_report.tex
%   pdflatex financial_sde_report.tex
% The experiment artifacts are regenerated from the repository root with:
%   python financial_sde_project.py
\documentclass[11pt,a4paper]{article}

\usepackage[utf8]{inputenc}
\usepackage[T1]{fontenc}
\usepackage{amsmath,amssymb,amsthm,bm,mathtools}
\usepackage[a4paper,margin=2.35cm]{geometry}
\usepackage{graphicx}
\graphicspath{{generated/}}
\usepackage{booktabs,array,longtable,multirow}
\usepackage{listings}
\usepackage{xcolor}
\usepackage{hyperref}
\usepackage{enumitem}
\usepackage{float}
\usepackage{caption}
\usepackage{subcaption}
\usepackage{algorithm}
\usepackage{algpseudocode}
\usepackage{mdframed}
\usepackage{microtype}
\usepackage{mathptmx}
\usepackage{url}

\definecolor{bluecb}{RGB}{0,119,187}
\definecolor{orangecb}{RGB}{238,119,51}
\definecolor{formbg}{RGB}{250,250,245}
\definecolor{formframe}{RGB}{180,170,140}
\hypersetup{colorlinks=true,linkcolor=bluecb,citecolor=bluecb,urlcolor=bluecb}
\setlength{\parindent}{1.5em}
\setlength{\parskip}{0.25em}
\captionsetup{font=small,labelfont=bf}

\lstset{
  language=Python,
  basicstyle=\ttfamily\scriptsize,
  keywordstyle=\color{blue}\bfseries,
  commentstyle=\color{green!40!black},
  stringstyle=\color{red!65!black},
  numbers=left,
  numberstyle=\tiny\color{gray},
  frame=single,
  breaklines=true,
  showstringspaces=false,
  tabsize=4,
  captionpos=b
}

\newtheorem{theorem}{Theorem}[section]
\newtheorem{proposition}[theorem]{Proposition}
\theoremstyle{definition}
\newtheorem{definition}[theorem]{Definition}
\newtheorem{assumption}[theorem]{Assumption}
\theoremstyle{remark}
\newtheorem{remark}[theorem]{Remark}

\newmdenv[
  backgroundcolor=formbg,linecolor=formframe,linewidth=0.8pt,
  roundcorner=3pt,innermargin=6pt,outermargin=5pt,
  innertopmargin=7pt,innerbottommargin=7pt
]{formbox}

\algrenewcommand{\algorithmicrequire}{\textbf{Input:}}
\algrenewcommand{\algorithmicensure}{\textbf{Output:}}

\title{\textbf{Reliable Numerical Simulation of Financial SDEs:}\\
\textbf{GBM Benchmarks and a Non-Affine Volatility Model}\\[2mm]
\large Project 1 (ODE IVP Extension Pathway: Financial SDE)}
\author{
  Team \#\texttt{<TEAM NUMBER>}\\
  \small \texttt{<Name 1>, <ID 1> (Project Manager)} \quad
  \texttt{<Name 2>, <ID 2> (Mathematical Theory)}\\
  \small \texttt{<Name 3>, <ID 3> (Algorithm Implementation)} \quad
  \texttt{<Name 4>, <ID 4> (Visualization \& Report)}\\
  \small \texttt{<Name 5>, <ID 5> (Testing \& Validation)}\\
  \small Department of Applied Mathematics
}
\date{18 September 2026}

\begin{document}
\maketitle

\begin{abstract}
This project develops and validates a reproducible numerical workflow for two
financial stochastic differential equations (SDEs).  Geometric Brownian motion
(GBM) is first used as an exact benchmark: Euler--Maruyama (EM), scalar Milstein,
and exact lognormal transitions are driven by identical Brownian increments.
This coupling isolates time-discretisation error from sampling noise.  With
$300{,}000$ paths, the fitted terminal strong-convergence slopes are
$0.4899\pm0.0026$ for EM and $0.9099\pm0.0167$ for Milstein, consistent with the
theoretical orders $1/2$ and $1$.  The exactly evaluable bias in
$\mathbb{E}[S_T]$ has fitted slope $0.9982\pm0.0004$.  At a matched mean absolute
error target of $0.85$ price units, EM uses 128 path steps while Milstein uses
2 under the declared step-count cost measure.  The experiment also checks the
terminal lognormal law and demonstrates that price-space updates are not
unconditionally positivity preserving.  We then solve a bounded, non-affine
stochastic-volatility model in log price.  Coarse paths are coupled to an
$N_{\mathrm{ref}}=2048$ EM reference by summing fine-grid Brownian increments.
The empirical strong slope is $0.5111\pm0.0047$; Brownian variance and
correlation diagnostics pass, and all simulated log-price paths remain finite
and positive.  Because the reference is numerical and call-payoff bias
confidence intervals include zero, no unsupported weak-order claim is made.
\end{abstract}

\noindent\textbf{Keywords:} stochastic differential equations; geometric
Brownian motion; Euler--Maruyama; Milstein method; strong convergence; weak
convergence; stochastic volatility; Monte Carlo coupling

\tableofcontents
\newpage

% ============================================================
\section{Introduction and motivation}
% ============================================================

\subsection{Problem background}
An SDE model is useful only if its numerical error can be separated from Monte
Carlo uncertainty.  GBM is therefore more than a toy model: its exact pathwise
transition, moments, and terminal distribution provide an oracle against which
an implementation can be tested.  Higham's algorithmic treatment emphasises
precisely this role of coupled paths when illustrating EM and Milstein
\cite{higham2001}; the broader convergence theory is developed by Kloeden and
Platen \cite{kloedenplaten1992} and Milstein \cite{milstein1995}.

The project asks a second, more practically relevant question: what changes
when no exact transition is available?  We answer it with the prescribed
two-factor non-affine stochastic-volatility model.  The same Brownian path is
represented on nested grids, so agreement under refinement becomes evidence
for convergence.  It is not relabelled as an exact error.  This distinction is
central to reliable computational finance and underlies multilevel coupling
methods \cite{giles2008}.

\subsection{Research questions and contribution}
The numerical investigation addresses four questions.
\begin{enumerate}[leftmargin=2.2em]
  \item Do coupled simulations recover the predicted GBM strong orders
  $1/2$ (EM) and $1$ (Milstein), and weak order $1$ for the mean?
  \item At a matched strong-error target, which method uses fewer time-step
  coefficient evaluations?
  \item When can price-space updates lose positivity, and does the simulated
  terminal law agree with the exact lognormal distribution?
  \item For the non-affine model, do nested-grid EM solutions converge, pass
  Brownian-increment checks, remain finite and positive, and use a sufficiently
  resolved numerical reference?
\end{enumerate}

The main contribution is a single, auditable script that implements every
scheme, fixes all seeds, regenerates every figure and printed number, and runs
explicit PASS/FAIL checks.  The analysis deliberately separates theorem,
measurement, and limitation.

\subsection{Report map}
Section~\ref{sec:model} defines both models.  Section~\ref{sec:methods} derives
the numerical schemes and convergence estimators.  Section~\ref{sec:implementation}
documents the implementation and validation design.  Section~\ref{sec:results}
presents GBM and non-affine experiments.  Section~\ref{sec:limitations}
identifies what the evidence does not establish, and Section~\ref{sec:conclusion}
concludes.  The complete source and AI/individual-contribution records appear
in the appendices.

% ============================================================
\section{Financial SDE models}\label{sec:model}
% ============================================================

\subsection{Geometric Brownian motion}
Let $(W_t)_{t\ge0}$ be a standard Brownian motion.  Under the physical measure,
GBM is
\begin{equation}
  dS_t=\mu S_t\,dt+\sigma S_t\,dW_t,\qquad S_0>0.
  \label{eq:gbm}
\end{equation}
For $f(s)=\log s$, It\^o's formula gives
\begin{align}
  d\log S_t
  &=f'(S_t)dS_t+\tfrac12 f''(S_t)(dS_t)^2 \notag\\
  &=\left(\mu-\tfrac12\sigma^2\right)dt+\sigma dW_t.
  \label{eq:ito-log}
\end{align}
Integrating~\eqref{eq:ito-log} yields the exact solution
\begin{equation}
  S_t=S_0\exp\!\left[\left(\mu-\tfrac12\sigma^2\right)t+\sigma W_t\right].
  \label{eq:gbm-exact}
\end{equation}
Consequently,
\begin{equation}
  \mathbb E[S_T]=S_0e^{\mu T},\qquad
  \operatorname{Var}(S_T)=S_0^2e^{2\mu T}\bigl(e^{\sigma^2T}-1\bigr),
  \label{eq:gbm-moments}
\end{equation}
and $S_T$ is lognormal.  The reproducible benchmark uses
\begin{equation}
  S_0=100,\qquad \mu=0.05,\qquad \sigma=0.40,\qquad T=1.
  \label{eq:gbm-params}
\end{equation}
The volatility is intentionally high enough to make the coarse-step positivity
question visible while keeping a familiar annual horizon.

\subsection{Non-affine stochastic volatility}
Write $X_t=\log S_t$.  The prescribed model is
\begin{align}
 dX_t &= \left(\mu-\tfrac12g(Y_t)^2\right)dt+g(Y_t)dW_t^{(1)},
 \label{eq:nonaff-x}\\
 dY_t &= \kappa(\theta-Y_t)dt+\xi\sqrt{1+Y_t^2}\,dW_t^{(2)},
 \label{eq:nonaff-y}\\
 g(y)&=\sigma_{\min}+\frac{\sigma_{\max}-\sigma_{\min}}{1+e^{-y}},
 \qquad d\langle W^{(1)},W^{(2)}\rangle_t=\rho\,dt.
 \label{eq:logistic}
\end{align}
The benchmark parameters are
\begin{equation}
\begin{gathered}
S_0=100,\quad Y_0=0,\quad \mu=0.05,\quad \kappa=2,
\quad\theta=-0.2,\quad\xi=0.6,\quad\rho=-0.7,\\
\sigma_{\min}=0.10,\quad\sigma_{\max}=0.50,\quad T=1.
\end{gathered}
\label{eq:nonaff-params}
\end{equation}
The logistic map guarantees $0.10<g(y)<0.50$.  Its derivative is bounded,
and $y\mapsto\sqrt{1+y^2}$ is globally Lipschitz because the magnitude of its
derivative is at most one.  Thus the coefficients have the regularity and
linear-growth structure needed for standard EM strong-convergence arguments;
nevertheless, the finite reference used below remains a numerical benchmark,
not a closed-form oracle.

% ============================================================
\section{Discretisation and mathematical analysis}\label{sec:methods}
% ============================================================

\subsection{Euler--Maruyama for GBM}
On $t_n=nh$, $h=T/N$, let $\Delta W_n=W_{t_{n+1}}-W_{t_n}sim
\mathcal N(0,h)$.  Freezing both coefficients at $S_n$ and integrating over one
step gives
\begin{equation}
  S_{n+1}^{\mathrm{EM}}
  =S_n^{\mathrm{EM}}+\mu S_n^{\mathrm{EM}}h
   +\sigma S_n^{\mathrm{EM}}\Delta W_n
  =S_n^{\mathrm{EM}}(1+\mu h+\sigma\Delta W_n).
  \label{eq:em-gbm}
\end{equation}
For globally Lipschitz drift and diffusion, EM has strong order $1/2$ and weak
order $1$ for sufficiently smooth test functions \cite{higham2001,kloedenplaten1992}.
GBM coefficients are globally Lipschitz, so the standard result applies.

\subsection{Milstein for GBM}
The scalar It\^o--Taylor expansion retains the iterated stochastic integral
\begin{equation}
 \int_{t_n}^{t_{n+1}}\int_{t_n}^{s}dW_r\,dW_s
 =\tfrac12\left[(\Delta W_n)^2-h\right].
 \label{eq:ito-iterated}
\end{equation}
For diffusion $b(s)=\sigma s$, $b'(s)=\sigma$, so the Milstein update is
\begin{equation}
 S_{n+1}^{\mathrm{Mil}}
 =S_n^{\mathrm{Mil}}+\mu S_n^{\mathrm{Mil}}h
 +\sigma S_n^{\mathrm{Mil}}\Delta W_n
 +\tfrac12\sigma^2S_n^{\mathrm{Mil}}\bigl[(\Delta W_n)^2-h\bigr].
 \label{eq:milstein-gbm}
\end{equation}
Including~\eqref{eq:ito-iterated} cancels the leading stochastic Taylor term
omitted by EM, raising the scalar strong order to $1$ while retaining weak order
$1$ \cite{milstein1995}.

\subsection{Exact transition and pathwise coupling}
For the same increment used in~\eqref{eq:em-gbm}--\eqref{eq:milstein-gbm},
\begin{equation}
 S_{n+1}^{\mathrm{exact}}
 =S_n^{\mathrm{exact}}
 \exp\!\left[\left(\mu-\tfrac12\sigma^2\right)h+\sigma\Delta W_n\right].
 \label{eq:exact-step}
\end{equation}
The exact, EM, and Milstein paths therefore share all $\Delta W_n$.  The terminal
strong-error estimator is
\begin{equation}
 \widehat e_{\mathrm{s}}(h)
 =\frac1M\sum_{m=1}^{M}\left|S_{N,m}^{(h)}-S_{T,m}^{\mathrm{exact}}\right|,
 \label{eq:strong-estimator}
\end{equation}
and a log--log least-squares slope estimates $p$ in
$\mathbb E|S_N-S_T|=O(h^p)$.  Pairing is essential: independent exact paths
would add $O(M^{-1/2})$ sampling variability unrelated to discretisation.

\subsection{Weak error and an analytic bias oracle}
For a test function $\varphi$, weak error is
\begin{equation}
 e_{\mathrm{w}}(h)=\left|\mathbb E\varphi(S_N^{(h)})-
 \mathbb E\varphi(S_T)\right|.
 \label{eq:weak-def}
\end{equation}
For $\varphi(s)=s$, both EM and Milstein step multipliers have conditional mean
$1+\mu h$.  Hence
\begin{equation}
 \mathbb E[S_N^{\mathrm{EM}}]
 =\mathbb E[S_N^{\mathrm{Mil}}]
 =S_0(1+\mu h)^N.
 \label{eq:numerical-mean}
\end{equation}
Using $N=T/h$ and $\log(1+\mu h)=\mu h-\tfrac12\mu^2h^2+O(h^3)$,
\begin{align}
 S_0(1+\mu h)^{T/h}
 &=S_0\exp\!\left(\mu T-\tfrac12\mu^2Th+O(h^2)\right),\\
 e_{\mathrm w}(h)
 &=\tfrac12S_0e^{\mu T}\mu^2T\,h+O(h^2).
 \label{eq:weak-expansion}
\end{align}
Thus the mean bias has exact asymptotic order one.  The code reports both this
deterministic bias and a paired Monte Carlo confidence interval for
$S_N-S_T$; the former prevents sampling noise from hiding the weak rate.

\subsection{Positivity analysis}
EM is positive for a step only when
\begin{equation}
 1+\mu h+\sigma\sqrt h Z_n>0,
\end{equation}
so the one-step failure probability is
\begin{equation}
 p_{\mathrm{neg,EM}}(h)
 =\Phi\!\left(-\frac{1+\mu h}{\sigma\sqrt h}\right)>0.
 \label{eq:em-neg-prob}
\end{equation}
For Milstein, the multiplier is the quadratic
\begin{equation}
 q(z)=\tfrac12\sigma^2hz^2+\sigma\sqrt h z
 +1+\mu h-\tfrac12\sigma^2h.
 \label{eq:milstein-quadratic}
\end{equation}
Its minimum is
\begin{equation}
 q_{\min}=\tfrac12+\mu h-\tfrac12\sigma^2h.
 \label{eq:milstein-min}
\end{equation}
For~\eqref{eq:gbm-params} and $0<h\le1$, this is positive, explaining why no
Milstein positivity violation is observed here.  Equation~\eqref{eq:milstein-min}
also shows that this is parameter- and step-dependent, not unconditional
positivity preservation.  Exact/log-price simulation remains positive by
construction.

\subsection{Correlated EM for the non-affine model}
Generate independent $Z_{1,n},Z_{2,n}\sim\mathcal N(0,1)$ and set
\begin{align}
 \Delta W_n^{(1)}&=\sqrt h Z_{1,n},\\
 \Delta W_n^{(2)}&=\sqrt h\left(\rho Z_{1,n}+\sqrt{1-\rho^2}Z_{2,n}\right).
 \label{eq:corr-inc}
\end{align}
The EM update is
\begin{align}
 X_{n+1}&=X_n+\left(\mu-\tfrac12g(Y_n)^2\right)h
              +g(Y_n)\Delta W_n^{(1)},
 \label{eq:em-x}\\
 Y_{n+1}&=Y_n+\kappa(\theta-Y_n)h
              +\xi\sqrt{1+Y_n^2}\Delta W_n^{(2)},
 \label{eq:em-y}\\
 S_{n+1}&=e^{X_{n+1}}>0.
 \label{eq:recover-s}
\end{align}
The logistic function is evaluated with two algebraically equivalent branches
to avoid overflow when $|Y_n|$ is large.

A scalar Milstein formula is \emph{not} applied to
\eqref{eq:nonaff-x}--\eqref{eq:nonaff-y}.  Written in two independent Brownian
drivers, the diffusion vector fields have nonzero cross derivatives through
$g(Y)$ and $\sqrt{1+Y^2}$.  A full multidimensional Milstein discretisation
would generally require cross terms and simulated iterated stochastic
integrals.  Omitting them while claiming strong order one would be
mathematically unjustified.

\subsection{Nested-grid numerical reference}
Let $h_{\mathrm{ref}}=T/N_{\mathrm{ref}}$.  For a coarse step $h=Kh_{\mathrm{ref}}$,
the coupling is
\begin{equation}
 \Delta W_{n,\mathrm{coarse}}^{(i)}
 =\sum_{j=0}^{K-1}\Delta W_{Kn+j,\mathrm{ref}}^{(i)},\qquad i=1,2.
 \label{eq:nested-sum}
\end{equation}
Therefore coarse and fine solutions represent the same Brownian path.  We use
$N_{\mathrm{ref}}=2048$ for the main experiment and independently compare it
with 4096 steps on 4000 fresh paths.  For the call payoff
$P(S_T)=(S_T-100)^+$, paired differences estimate weak bias and substantially
reduce variance.  We report a weak trend only where its confidence interval is
separated from zero.

\subsection{Uncertainty quantification}
For independent path-level observations $D_1,\ldots,D_M$, every displayed
Monte Carlo mean uses
\begin{equation}
 \overline D\ \pm\ 1.96\frac{s_D}{\sqrt M},
 \qquad
 s_D^2=\frac{1}{M-1}\sum_{m=1}^{M}(D_m-\overline D)^2.
 \label{eq:ci}
\end{equation}
The bands around fitted log--log lines are regression diagnostics, not repeated
Monte Carlo uncertainty.  Their slope standard errors quantify linear-fit
uncertainty conditional on the chosen mesh range.

% ============================================================
\section{Implementation and reproducibility}\label{sec:implementation}
% ============================================================

\subsection{Core algorithm}
Algorithm~\ref{alg:nested} states the non-affine nested-grid calculation.  GBM
uses the same coupling pattern with the exact step~\eqref{eq:exact-step}.

\begin{algorithm}[H]
\caption{Coupled non-affine EM convergence experiment}
\label{alg:nested}
\begin{algorithmic}[1]
\Require Parameters~\eqref{eq:nonaff-params}, seed, path count $M$, fine count
$N_{\rm ref}$, coarse counts $\mathcal N$
\Ensure Strong errors, paired payoff biases, confidence intervals, diagnostics
\For{each memory-safe batch of paths}
  \State draw independent fine-grid arrays $Z_1,Z_2$
  \State construct correlated increments using~\eqref{eq:corr-inc}
  \State integrate~\eqref{eq:em-x}--\eqref{eq:em-y} on the reference grid
  \For{$N\in\mathcal N$}
    \State sum groups of $N_{\rm ref}/N$ increments using~\eqref{eq:nested-sum}
    \State integrate the coarse path and store paired terminal differences
  \EndFor
\EndFor
\State estimate means and confidence intervals using~\eqref{eq:ci}
\State validate increment variance/correlation, finiteness, and positivity
\end{algorithmic}
\end{algorithm}

\subsection{Single-file code structure}
The requested submission has two primary source files.  Accordingly, the code
is organised as small, model-agnostic functions inside
\texttt{financial\_sde\_project.py}, rather than split across modules.

\begin{table}[H]
\centering
\small
\begin{tabular}{p{0.34\textwidth}p{0.58\textwidth}}
\toprule
\textbf{Function group} & \textbf{Purpose}\\
\midrule
\texttt{gbm\_terminal\_coupled}, \texttt{run\_gbm}
& Exact/EM/Milstein coupling, strong and weak estimators, positivity and matched-cost results.\\
\texttt{logistic\_vol}, \texttt{em\_nonaffine\_from\_increments}
& Stable model coefficients and swappable EM integration from supplied increments.\\
\texttt{aggregate\_increments}, \texttt{run\_nonaffine}
& Nested Brownian coupling, reference comparison, call payoff and diagnostics.\\
\texttt{nonaffine\_reference\_refinement}
& Independent $2048\rightarrow4096$ reference-resolution check.\\
\texttt{plot\_*}, \texttt{validate}
& Reproduce all five figures and execute ten PASS/FAIL checks.\\
\bottomrule
\end{tabular}
\caption{Code organisation.  Integrators receive increments explicitly, so the
random driver and numerical method can be tested independently.}
\label{tab:code-structure}
\end{table}

\subsection{Reproducible protocol}
The default run uses seed 20260918, $M=300{,}000$ GBM paths, $M=20{,}000$
non-affine paths, batches of at most 1000 paths, and $N_{\rm ref}=2048$.  The
command
\begin{center}
\texttt{python financial\_sde\_project.py}
\end{center}
regenerates \texttt{latex/generated/results.json},
\texttt{run\_output.txt}, and Figures~\ref{fig:gbm-convergence}--\ref{fig:nonaff-paths}
with fixed names and at 220 DPI.  No absolute paths are used.  A reduced
\texttt{--quick} mode is available only as a smoke test; all report numbers are
from the default run.  NumPy performs vectorised path operations
\cite{harris2020}, and Matplotlib produces the figures \cite{hunter2007}.

\subsection{Oracle chain and automatic checks}
Two independent GBM oracles are used: the pathwise exact transition and the
closed-form terminal moments/distribution.  The non-affine reference is checked
by an independent finer grid.  Ten automatic checks cover expected-order
bands, the exact GBM mean, both increment variances, target correlation,
finiteness, positivity, and a plausible non-affine strong slope.  All ten pass
in the archived output.  The checks are deliberately broad enough to avoid
passing only one random realisation, yet narrow enough to catch a wrong
increment scaling, correlation formula, or scheme.

\subsection{Declared cost measure}
The cost measure is the number of time steps per path, equivalently one
drift/diffusion coefficient evaluation per step in the implemented scalar GBM
loop.  The comparison excludes plotting, random-number setup, file I/O,
compilation, and vector-allocation overhead.  It also does not pretend that an
EM and a Milstein step have identical floating-point operation counts: Milstein
evaluates an extra correction.  The reported ratio is therefore an algorithmic
step-count comparison, not a wall-clock speedup claim.

% ============================================================
\section{Numerical experiments and results}\label{sec:results}
% ============================================================

\subsection{GBM strong and weak convergence}
Table~\ref{tab:gbm-convergence} reports the asymptotic mesh range used for the
fit.  Figure~\ref{fig:gbm-convergence} overlays the predicted slopes and 95\%
fit bands.  The observed EM slope $0.4899\pm0.0026$ agrees with $1/2$; the
Milstein slope $0.9099\pm0.0167$ approaches order one.  The small shortfall from
one is consistent with a finite pre-asymptotic mesh range rather than evidence
for a different theoretical order.  The analytic mean-bias slope
$0.9982\pm0.0004$ is essentially one.

\begin{table}[H]
\centering
\small
\begin{tabular}{rrrrr}
\toprule
$N$ & $h$ & EM strong MAE & Milstein strong MAE & exact mean bias magnitude\\
\midrule
4   & 0.250000 & 4.52887 & 0.460308 & 0.0325759\\
8   & 0.125000 & 3.28870 & 0.268833 & 0.0163567\\
16  & 0.062500 & 2.34715 & 0.149881 & 0.00819567\\
32  & 0.031250 & 1.67424 & 0.0800666 & 0.00410218\\
64  & 0.015625 & 1.18565 & 0.0414826 & 0.00205218\\
128 & 0.0078125 & 0.838329 & 0.0211342 & 0.00102636\\
256 & 0.00390625 & 0.594648 & 0.0106906 & 0.000513248\\
\bottomrule
\end{tabular}
\caption{GBM convergence with $300{,}000$ coupled paths.  Error decreases at
the theoretically expected rates; the exact weak-bias column is common to EM
and Milstein for the linear payoff $\varphi(s)=s$.}
\label{tab:gbm-convergence}
\end{table}

\begin{figure}[H]
\centering
\includegraphics[width=0.98\textwidth]{fig01_gbm_convergence.png}
\caption{Coupled GBM convergence.  Left: EM and Milstein terminal mean absolute
errors follow the $h^{1/2}$ and $h$ reference trends, respectively; fit bands
make the slope uncertainty visible.  Right: the analytically evaluable mean
bias follows order one.  This separates discretisation error from Monte Carlo
noise and directly supports the claimed orders.}
\label{fig:gbm-convergence}
\end{figure}

\subsection{Matched-accuracy cost comparison}
At the predeclared target
$\mathbb E|S_N-S_T|\le0.85$, the coarsest tested EM grid that passes uses
$N=128$ and attains 0.83833; Milstein passes at $N=2$ with 0.75433.
Under the stated step-count measure, the work ratio is therefore
\begin{equation}
 \frac{C_{\rm EM}}{C_{\rm Mil}}=\frac{128}{2}=64.
\end{equation}
This is a valid matched-accuracy comparison, but it is not a claim of a
64-fold wall-clock speedup.  The negative result implicit in the comparison is
also important: at very loose accuracy, EM's cheaper algebra could matter;
that regime is not measured by this target.

\subsection{Terminal distribution and moment validation}
The exact formulas give
\begin{equation}
 \mathbb E[S_T]=105.127110,qquad
 \operatorname{Var}(S_T)=1917.591686.
\end{equation}
The exact-transition Monte Carlo sample gives 105.088894 and 1915.728882,
respectively; the mean is within four Monte Carlo standard errors, satisfying
the automatic oracle check.  At $N=64$, the 1\%, 5\%, 50\%, 95\%, and 99\%
quantiles are shown in Table~\ref{tab:quantiles}.

\begin{table}[H]
\centering
\small
\begin{tabular}{lrrrrr}
\toprule
Method & 1\% & 5\% & 50\% & 95\% & 99\%\\
\midrule
Exact lognormal & 38.2687 & 50.2610 & 97.0446 & 187.3749 & 246.0930\\
EM, $N=64$ & 37.9146 & 50.0592 & 97.0464 & 187.5114 & 244.2812\\
Milstein, $N=64$ & 38.3043 & 50.2674 & 96.9520 & 187.6476 & 245.1338\\
\bottomrule
\end{tabular}
\caption{GBM terminal quantiles in price units.  Central quantiles agree
closely; the tails reveal the remaining finite-step error more sensitively.}
\label{tab:quantiles}
\end{table}

\begin{figure}[H]
\centering
\includegraphics[width=0.98\textwidth]{fig02_gbm_distribution.png}
\caption{Terminal-law validation at $N=64$.  The density overlay checks the
whole distribution rather than only two moments, while the quantile plot makes
tail discrepancies visible.  Milstein tracks the exact lognormal quantiles
more closely than EM at the same step count.}
\label{fig:gbm-distribution}
\end{figure}

\subsection{Positivity stress test}
At $h=1$, 0.435\% of EM paths encounter a non-positive step multiplier; at
$h=1/2$, the rate is 0.034\%, and no violation is observed among 300,000 paths
for $h\le1/4$.  This is consistent with the rapid Gaussian-tail decay in
\eqref{eq:em-neg-prob}.  No Milstein violation occurs for the tested parameters,
as predicted by~\eqref{eq:milstein-min}.

\begin{figure}[H]
\centering
\includegraphics[width=0.72\textwidth]{fig03_gbm_positivity.png}
\caption{Observed fraction of paths with at least one non-positive update
multiplier.  EM fails on coarse grids even though the true GBM remains positive;
zeros at finer grids mean ``none in 300,000 paths'', not a mathematical
guarantee.  Milstein is positive here because its quadratic multiplier has a
positive minimum for the tested $h$, not because Milstein is unconditionally
positivity preserving.}
\label{fig:gbm-positivity}
\end{figure}

\subsection{Non-affine nested-grid convergence}
Table~\ref{tab:nonaffine} compares each coarse grid with the coupled 2048-step
reference.  Fitting the first four rows (the final row is reserved as a
near-reference diagnostic) gives strong slope $0.5111\pm0.0047$, closely
matching the EM benchmark $1/2$.  Every 95\% confidence interval for the paired
call-payoff difference contains zero.  The correct conclusion is therefore
that sampling uncertainty is already comparable to the estimated weak bias;
the data do not support quoting a weak slope.

\begin{table}[H]
\centering
\small
\begin{tabular}{rrrrr}
\toprule
$N$ & $h$ & strong MAE & strong 95\% CI & call bias (95\% CI)\\
\midrule
16  & 0.062500 & 0.886321 & [0.875404, 0.897238] & 0.004204 $[-0.010095,0.018504]$\\
32  & 0.031250 & 0.627466 & [0.619890, 0.635042] & 0.003974 $[-0.006056,0.014005]$\\
64  & 0.015625 & 0.441095 & [0.435856, 0.446335] & 0.003689 $[-0.003286,0.010664]$\\
128 & 0.0078125 & 0.306001 & [0.302342, 0.309660] & 0.002141 $[-0.002714,0.006996]$\\
256 & 0.00390625 & 0.210677 & [0.208183, 0.213171] & 0.002748 $[-0.000568,0.006064]$\\
\bottomrule
\end{tabular}
\caption{Non-affine EM convergence against the coupled $N_{\rm ref}=2048$
numerical reference using 20,000 paths.  Strong convergence is resolved, while
the weak payoff bias is not statistically distinguishable from zero.}
\label{tab:nonaffine}
\end{table}

\begin{figure}[H]
\centering
\includegraphics[width=0.98\textwidth]{fig04_nonaffine_convergence.png}
\caption{Nested-grid evidence for the non-affine model.  Left: the pathwise
terminal error follows an empirical half-order trend with narrow Monte Carlo
intervals.  Right: all payoff-bias confidence intervals meet zero, so claiming
a weak order from these samples would over-interpret noise.}
\label{fig:nonaff-convergence}
\end{figure}

\subsection{Increment, finiteness, positivity, and reference checks}
The fine-grid increment diagnostics are
\begin{equation}
 \frac{\widehat{\operatorname{Var}}(\Delta W^{(1)})}{h}=1.00526,
 \qquad
 \frac{\widehat{\operatorname{Var}}(\Delta W^{(2)})}{h}=1.00346,
 \qquad
 \widehat{\operatorname{Corr}}=-0.70185,
\end{equation}
against targets $(1,1,-0.7)$.  All 20,000 reference paths are finite and
positive.  Their estimated mean terminal price is $105.01815$ with standard
error $0.21433$; the call value estimate is $14.48371$ with standard error
$0.15090$.  These are model outputs, not closed-form benchmarks.

On 4000 fresh coupled paths, refining the numerical reference from 2048 to 4096
steps changes the terminal price by mean absolute amount
\begin{equation}
 0.057626\quad\text{with 95\% CI }[0.056114,0.059139].
 \label{eq:refinement-result}
\end{equation}
This is substantially smaller than the coarsest-grid strong errors and about
27\% of the $N=256$ error, so the reference is adequate for the observed trend,
though not exact.  A still finer reference would be required before reporting
more terminal digits or resolving the weak payoff bias.

\begin{figure}[H]
\centering
\includegraphics[width=0.98\textwidth]{fig05_nonaffine_paths.png}
\caption{Representative non-affine paths.  Integrating $X=\log S$ guarantees
positive displayed prices after exponentiation (left), while the logistic
map keeps instantaneous volatility strictly between 0.10 and 0.50
year$^{-1/2}$ (right).  These structural properties are enforced by the model
representation rather than inferred from a finite sample.}
\label{fig:nonaff-paths}
\end{figure}

% ============================================================
\section{Limitations and robustness}\label{sec:limitations}
% ============================================================
\begin{enumerate}[leftmargin=2.2em]
  \item \textbf{Numerical reference.}  The non-affine error is relative to EM
  with 2048 steps.  Equation~\eqref{eq:refinement-result} quantifies, but does
  not eliminate, reference bias.
  \item \textbf{Weak convergence resolution.}  Paired common random numbers
  reduce variance, yet the call-bias intervals still contain zero.  The report
  therefore does not manufacture a weak slope.  More paths, multilevel Monte
  Carlo, or a different payoff would be needed.
  \item \textbf{Regression uncertainty.}  Slope standard errors diagnose the
  chosen log--log fit.  They are not theorem proofs and do not include every
  model-selection uncertainty, such as mesh-range choice.
  \item \textbf{Cost definition.}  Step count is portable but omits the extra
  arithmetic in Milstein and hardware effects.  A wall-clock study should use
  repeated timings on the same machine and separately count random draws and
  coefficient/derivative evaluations.
  \item \textbf{Model risk.}  Numerical convergence does not validate GBM or
  the non-affine model as a description of market data.  Calibration,
  risk-neutral drift, transaction costs, jumps, and empirical fit are outside
  this numerical-analysis project.
  \item \textbf{Extreme paths.}  The stable logistic evaluation and batched
  finiteness checks reduce numerical risk, but $e^X$ can overflow under much
  longer horizons or more extreme parameters.  Production code should retain
  log values and define an overflow policy.
\end{enumerate}

% ============================================================
\section{Conclusions}\label{sec:conclusion}
% ============================================================
\begin{itemize}[leftmargin=2em]
  \item Exact GBM coupling recovered the expected strong orders: EM
  $0.4899\pm0.0026$ and Milstein $0.9099\pm0.0167$; the mean weak bias had order
  $0.9982\pm0.0004$.
  \item At matched strong MAE $\le0.85$, Milstein required 2 steps versus 128
  for EM under the declared step-count cost, while the report avoids converting
  that 64:1 ratio into an unsupported wall-clock claim.
  \item Distribution and moment checks agree with the exact lognormal oracle,
  and the coarse-grid stress test shows why price-space EM is not
  positivity-preserving.
  \item For the non-affine model, nested Brownian coupling produced empirical
  half-order strong convergence and passed variance, correlation, finiteness,
  positivity, and reference-refinement checks.  Weak payoff bias remained
  unresolved at the available Monte Carlo precision.
  \item A natural extension is multilevel Monte Carlo, which would reuse the
  same nested coupling to reduce the cost of resolving weak payoffs
  \cite{giles2008}; a second extension is a correctly implemented
  multidimensional Milstein scheme with iterated stochastic integrals.
\end{itemize}

% ============================================================
\section*{References}
\addcontentsline{toc}{section}{References}
% ============================================================
\begin{thebibliography}{99}
\bibitem{higham2001}
D. J. Higham, ``An algorithmic introduction to numerical simulation of
stochastic differential equations,'' \emph{SIAM Review}, vol.~43, no.~3,
pp.~525--546, 2001. doi: \href{https://doi.org/10.1137/S0036144500378302}{10.1137/S0036144500378302}.

\bibitem{kloedenplaten1992}
P. E. Kloeden and E. Platen, \emph{Numerical Solution of Stochastic
Differential Equations}. Berlin: Springer, 1992.
doi: \href{https://doi.org/10.1007/978-3-662-12616-5}{10.1007/978-3-662-12616-5}.

\bibitem{milstein1995}
G. N. Milstein, \emph{Numerical Integration of Stochastic Differential
Equations}. Dordrecht: Springer, 1995.
doi: \href{https://doi.org/10.1007/978-94-015-8455-5}{10.1007/978-94-015-8455-5}.

\bibitem{giles2008}
M. B. Giles, ``Multilevel Monte Carlo path simulation,''
\emph{Operations Research}, vol.~56, no.~3, pp.~607--617, 2008.
doi: \href{https://doi.org/10.1287/opre.1070.0496}{10.1287/opre.1070.0496}.

\bibitem{harris2020}
C. R. Harris et al., ``Array programming with NumPy,'' \emph{Nature},
vol.~585, pp.~357--362, 2020.
doi: \href{https://doi.org/10.1038/s41586-020-2649-2}{10.1038/s41586-020-2649-2}.

\bibitem{hunter2007}
J. D. Hunter, ``Matplotlib: A 2D graphics environment,''
\emph{Computing in Science \& Engineering}, vol.~9, no.~3, pp.~90--95, 2007.
doi: \href{https://doi.org/10.1109/MCSE.2007.55}{10.1109/MCSE.2007.55}.

\bibitem{coursebrief}
MTH321 teaching team, \emph{Project 1 Brief: Numerical Solution of ODE Initial
Value Problems}, extension pathway 4, 2026, course material supplied with the
project.

\bibitem{problempack}
MTH321 teaching team, \emph{Problem Pack: Five Suggested Directions for the ODE
IVP Project}, ``Financial SDE: GBM benchmark and non-affine volatility,'' 2026,
course material supplied with the project.
\end{thebibliography}

\appendix

% ============================================================
\section{Reproduction record and complete console output}
% ============================================================
\subsection{Environment and commands}
Install the declared dependencies in any isolated Python environment:
\begin{center}
\texttt{python -m pip install numpy matplotlib}
\end{center}
From the repository root run
\begin{center}
\texttt{python financial\_sde\_project.py}.
\end{center}
All paths are relative; the output directory is created automatically.  The
JSON file records Python, NumPy, Matplotlib, platform, seed, sample sizes,
parameters, every table row, validation results, and wall time.

\subsection{Archived program output}
\lstinputlisting[language={},caption={Complete output from the full reproducibility run.}]{generated/run_output.txt}

% ============================================================
\section{Complete Python source code}
% ============================================================
The following listing is the exact code used to generate every reported result
and figure.
\lstinputlisting[caption={Complete source: \texttt{financial\_sde\_project.py}.}]{../financial_sde_project.py}

% ============================================================
\section{AI Transparency Log}
% ============================================================
\subsection{Tools used}
\begin{table}[H]
\centering
\small
\begin{tabular}{p{0.25\textwidth}p{0.20\textwidth}p{0.46\textwidth}}
\toprule
Tool & Version/date & Primary use\\
\midrule
OpenAI Codex & 18 September 2026 & Requirements mapping, derivation drafting,
Python implementation, experiment execution, result interpretation, and
LaTeX assembly.\\
NumPy / Matplotlib & Versions recorded in \texttt{results.json} & Vectorised
Monte Carlo computation and deterministic figure generation.\\
\bottomrule
\end{tabular}
\caption{AI and software tools.  Team members must replace or supplement this
record if they make later changes with other tools.}
\end{table}

\subsection{Significant AI-assisted claim audits}
\begin{longtable}{p{0.06\textwidth}p{0.43\textwidth}p{0.43\textwidth}}
\toprule
Item & Prompt/claim/assumptions & Risk and possible falsifier\\
\midrule
1 & Derive and test the GBM strong orders of EM and Milstein under pathwise
coupling. & A wrong increment scale or uncoupled oracle can create a false
slope.  Falsifiers: exact-transition errors fail to decrease or fitted slopes
leave the declared acceptance bands.\\
2 & Use a nested EM reference for the non-affine model and describe the strong
rate as empirical. & Treating a discretised reference as exact biases the
claim.  Falsifier: the 2048-to-4096 reference difference is comparable to the
coarse-grid error or refinement ceases to reduce error.\\
3 & Compare EM and Milstein cost at matched accuracy. & Counting only steps can
overstate a speed advantage because a Milstein step has extra algebra.
Falsifier: repeated wall-clock or operation-count measurements reverse the
comparison at the same error target.\\
\bottomrule
\end{longtable}

\subsection{Independent verification and decisions}
\begin{longtable}{p{0.06\textwidth}p{0.46\textwidth}p{0.40\textwidth}}
\toprule
Item & Independent check and evidence location & Final decision/modification/limitation\\
\midrule
1 & Exact GBM transition, exact moments, analytic mean bias, Tables
\ref{tab:gbm-convergence}--\ref{tab:quantiles}, and automatic order checks. &
Accepted the strong/weak claims; reported fitted standard errors and separated
analytic weak bias from Monte Carlo estimates.\\
2 & Fresh coupled 2048-to-4096 refinement, Eq.~\eqref{eq:refinement-result},
plus increment variance/correlation checks. & Called the result ``empirical'';
did not describe the reference as exact and did not claim a weak payoff slope.\\
3 & Same strong-MAE threshold for both schemes and explicit cost exclusions in
Section~\ref{sec:implementation}. & Reported a 64:1 step-count ratio only;
explicitly withheld a wall-clock speedup claim.\\
\bottomrule
\end{longtable}

\subsection{Intellectual ownership declaration}
Before submission, every member must complete and sign the table below.  Each
signer affirms that they can derive the equations, explain and modify the
team-owned code, identify trusted library interfaces, and reproduce the
validation without AI assistance during the viva.

\begin{table}[H]
\centering
\small
\begin{tabular}{p{0.17\textwidth}p{0.15\textwidth}p{0.35\textwidth}p{0.24\textwidth}}
\toprule
Name & Signature & Team-owned artefact/contribution & Trusted dependency/interface\\
\midrule
\texttt{<Name 1>} & & & \\
\texttt{<Name 2>} & & & \\
\texttt{<Name 3>} & & & \\
\texttt{<Name 4>} & & & \\
\texttt{<Name 5>} & & & \\
\bottomrule
\end{tabular}
\end{table}

% ============================================================
\section{Individual Contribution Statements}
% ============================================================
The following forms are intentionally not pre-filled: contribution claims must
be written by the corresponding student and remain traceable to the team's Git
history.  Copy the appropriate concrete evidence (function, equation, figure,
test, or report section) rather than writing ``worked on the report.''

\newcommand{\icsform}[2]{
\begin{formbox}
\textbf{Individual Contribution Statement --- Project 1 (Financial SDE)}\\[0.4em]
\textbf{Name:} \texttt{<#1>} \hfill \textbf{Role:} \texttt{<#2>}\\[0.4em]
\textbf{The specific things I did this project:}
\begin{enumerate}[leftmargin=2em]
\item \texttt{<Name a file/function/equation/figure and the exact work done.>}
\item \texttt{<Name a validation or review action and its evidence location.>}
\item \texttt{<Name a report section or experiment and the decision made.>}
\end{enumerate}
\textbf{The hardest conceptual obstacle and how I overcame it:}\\
\texttt{<Specific mathematical or algorithmic difficulty.>}\\[0.4em]
\textbf{One AI-assisted item I used and what I changed:}\\
\texttt{<Cross-reference AI Transparency Log item 1, 2, or 3.>}\\[0.4em]
\textbf{What I would do differently / what I still need to understand:}\\
\texttt{<Honest reflection.>}\\[0.6em]
\textbf{Signature:} \underline{\hspace{4cm}}\qquad
\textbf{Date:} \underline{\hspace{2.5cm}}
\end{formbox}
\newpage
}

\icsform{Name 1}{Project Manager}
\icsform{Name 2}{Mathematical Theory}
\icsform{Name 3}{Algorithm Implementation}
\icsform{Name 4}{Visualization \& Report}
\icsform{Name 5}{Testing \& Validation}

% ============================================================
\section{Submission evidence matrix}
% ============================================================
\begin{longtable}{p{0.33\textwidth}p{0.58\textwidth}}
\toprule
\textbf{Course requirement} & \textbf{Evidence in this submission}\\
\midrule
SDE-specific methods and theory & Sections~\ref{sec:model}--\ref{sec:methods}:
exact GBM, EM, Milstein, strong/weak definitions, positivity, and coupled
non-affine EM.\\
Independent oracle & Exact transition, exact moments/distribution, analytic
mean bias, plus non-affine reference refinement.\\
Observed order for at least two methods & Table~\ref{tab:gbm-convergence} and
Figure~\ref{fig:gbm-convergence}, with slopes and fit uncertainty.\\
Monte Carlo uncertainty & Path-level 95\% confidence intervals in
Table~\ref{tab:nonaffine} and explicit formula~\eqref{eq:ci}.\\
Matched-accuracy cost & Section~\ref{sec:results}: common MAE target 0.85,
declared step-count cost, exclusions, and limitation.\\
Distribution and sanity checks & Figure~\ref{fig:gbm-distribution},
Table~\ref{tab:quantiles}, moment oracle, increment variance/correlation,
finiteness, positivity.\\
Non-affine nested validation & Equation~\eqref{eq:nested-sum},
Table~\ref{tab:nonaffine}, Figure~\ref{fig:nonaff-convergence}, and
2048-to-4096 refinement.\\
Reproducible code and figures & One command regenerates five fixed-name,
220-DPI figures, JSON, and console output; fixed seed; no absolute paths.\\
Interpretive visualisations & Every figure has units, a colorblind-safe palette,
and a caption stating both the evidence and its significance.\\
Scientific limitations & Section~\ref{sec:limitations} separates numerical
convergence, reference error, sampling uncertainty, cost, and model risk.\\
AI log and ICS & Appendices C and D, with audit/falsifier records and five
member-specific forms.\\
\bottomrule
\end{longtable}

\end{document}
